\documentclass[11pt,a4paper]{article}
\usepackage[margin=2.2cm]{geometry}
\usepackage{booktabs}
\usepackage{hyperref}
\usepackage{enumitem}
\usepackage{xcolor}
\usepackage{graphicx}
\usepackage{subcaption}

\title{\textbf{OpenH-RF Proposal:}\\[4pt]
The UltraSound ToolBox Channel Capture Collection}
\author{
  Ole Marius Hoel Rindal,
  Y\"ucel Karabiyik,
  Sven Peter N\"asholm,
  Andreas Austeng
  \\[4pt]
  University of Oslo (UiO), Department of Informatics, Norway
}
\date{}

\begin{document}
\maketitle
\thispagestyle{empty}

% ------------------------------------------------------------------
\section{Executive Summary}

\subsection*{The UltraSound ToolBox (USTB)}
The USTB is a free, open-source MATLAB toolbox for beamforming, processing, and visualization of ultrasonic signals, first released in 2017~\cite{rodriguez2017ustb}.
It is developed at the University of Oslo in collaboration with NTNU, Universit\'e de Lyon, and Johns Hopkins University.
USTB provides a unified processing framework based on the Ultrasound File Format (UFF, HDF5-based) and a General Beamformer architecture supporting all common transmit sequences (plane wave, diverging wave, focused, synthetic aperture, and retrospective beamforming).
According to Google Scholar, the USTB has been \textbf{referenced in more than 150 publications}, spanning a wide range of topics such as adaptive beamforming, image quality assessment, deep learning for ultrasound, and clinical imaging research.

\subsection*{Proposed Contribution}
We propose to contribute \textbf{51 USTB channel-capture datasets to OpenH-RF}.
The collection is diverse in anatomy, hardware, and experimental and applied acquisition strategies:
\begin{itemize}[nosep]
  \item[\textbf{A}] \textbf{In-vivo human (cardiac)} --- parasternal long-axis and apical four-chamber imaging (3 acquisitions).
  \item[\textbf{B}] \textbf{In-vivo human (vascular)} --- carotid artery cross-section and longitudinal views (5 acquisitions).
  \item[\textbf{C}] \textbf{Phantoms (Verasonics)} --- tissue-mimicking CIRS targets for resolution, contrast, and dynamic range evaluation (18 acquisitions).
  \item[\textbf{D}] \textbf{Phantoms (Alpinion)} --- linear-array phantom benchmarks (4 acquisitions).
  \item[\textbf{E}] \textbf{Simulations} --- wave-propagation-based synthetic channel data (Field~II): point scatterers, cysts, speckle, and blocked-array configurations (17 acquisitions).
  \item[\textbf{F}] \textbf{Motion estimation} --- elastography phantoms for shear-wave and ARFI push-tracking (4 acquisitions).
\end{itemize}

\noindent All data is openly hosted on Zenodo (\url{https://zenodo.org/records/20261898}).
A catalog of all datasets, each with a reference Delay-And-Sum (DAS) beamformed B-mode image, is available at \url{https://unioslo.github.io/USTB/datasets.html}.

% ------------------------------------------------------------------
\section{Task \& Application Selection}

\subsection*{Hardware}
The experimental (non-simulated) channel data was recorded using two research ultrasound platforms, both providing full access to pre-beamformed RF channel data:
\begin{itemize}[nosep]
  \item \textbf{Verasonics Vantage 256} --- majority of datasets (cardiac, vascular, phantoms, SWE/ARFI)
  \item \textbf{Alpinion E-Cube 12R} --- phantom imaging with the L3-8 linear probe
\end{itemize}
\noindent Four transducer probes were used:
\begin{itemize}[nosep]
  \item \textbf{P4-2 phased array} (64 elements, Verasonics) --- cardiac and focused imaging
  \item \textbf{L7-4 linear array} (128 elements, Verasonics) --- vascular, phantom, and general imaging
  \item \textbf{L3-8 linear array} (128 elements, Alpinion) --- phantom imaging (4 datasets)
\end{itemize}

\subsection*{Transmit Strategies}
The collection covers five distinct transmit configurations:
coherent plane-wave compounding (CPWC, 5--75 angles),
diverging-wave imaging (DW, 25 virtual sources),
focused single-line imaging (FI, 128 focused beams),
synthetic transmit aperture (STA, 128 single-element transmits),
and shear-wave push-tracking sequences.

\subsection*{Dataset Breakdown}

\begin{table}[h!]
\centering
\small
\begin{tabular}{@{}cp{3.6cm}cp{1.8cm}cp{3.0cm}@{}}
\toprule
\textbf{ID} & \textbf{Application} & \textbf{No.\ of acq.} & \textbf{Tier} & \textbf{No.\ of frames} & \textbf{Probe / TX} \\
\midrule
\multicolumn{6}{@{}l}{\textit{6.1 Generalized Reconstruction}} \\
A & Echocardiography        & 3  & In-vivo (x40) & 80  & P4-2 / FI \\
B & Carotid imaging         & 5  & In-vivo (x40) & 8   & L7-4 / FI, CPWC \\
C & Phantom (Verasonics)    & 18 & Phantom (x1)  & 26  & L7-4, P4-2 / CPWC, FI, STA \\
D & Phantom (Alpinion)      & 4  & Phantom (x1)  & 4   & L3-8 / FI, CPWC \\
E & Simulated               & 17 & Simulation (x1) & 22 & Field~II / CPWC, STA, FI \\
\midrule
\multicolumn{6}{@{}l}{\textit{6.4 Motion Estimation}} \\
F & SWE / ARFI              & 4  & Phantom (x1)  & 800 & L7-4 / push-track \\
\midrule
& \textbf{Total}            & \textbf{51} & & \textbf{940} & \\
\bottomrule
\end{tabular}
\caption{Dataset breakdown by application group (labels A--F match the executive summary). No.\ of acq.\ = number of acquisitions; No.\ of frames = total channel-capture frames across all acquisitions in the group.}
\end{table}

\vspace{0.3em}
\noindent\textbf{Weighted contribution summary:}

\begin{center}
\small
\begin{tabular}{@{}lrrr@{}}
\toprule
\textbf{Tier} & \textbf{No.\ of datasets} & \textbf{Weight} & \textbf{Points} \\
\midrule
In-vivo human (research platform) & 8 & $\times 40$ & 320 \\
Phantom / table-top & 26 & $\times 1$ & 26 \\
Simulation (digital) & 17 & $\times 1$ & 17 \\
\midrule
\textbf{Total} & \textbf{51} & & \textbf{363} \\
\bottomrule
\end{tabular}
\end{center}
\noindent\textit{Weight denotes the OpenH-RF tier multiplier from the RFP (e.g.\ in-vivo human datasets on research platforms count $40\times$ toward the contribution total). Points $=$ No.\ of datasets $\times$ Weight.}

% ------------------------------------------------------------------
\begin{figure}[h!]
\centering
\small
\begin{subfigure}[t]{0.24\textwidth}
  \includegraphics[width=\textwidth]{thumbnails/cardiac_parasternal.png}
  \caption{Cardiac P4-2}
\end{subfigure}
\begin{subfigure}[t]{0.24\textwidth}
  \includegraphics[width=\textwidth]{thumbnails/cardiac_apical.png}
  \caption{Cardiac apical}
\end{subfigure}
\begin{subfigure}[t]{0.24\textwidth}
  \includegraphics[width=\textwidth]{thumbnails/carotid_cross.png}
  \caption{Carotid L7-4}
\end{subfigure}
\begin{subfigure}[t]{0.24\textwidth}
  \includegraphics[width=\textwidth]{thumbnails/carotid_long.png}
  \caption{Carotid long.}
\end{subfigure}

\vspace{0.3em}
\begin{subfigure}[t]{0.24\textwidth}
  \includegraphics[width=\textwidth]{thumbnails/cpwc_phantom.png}
  \caption{CPWC phantom}
\end{subfigure}
\begin{subfigure}[t]{0.24\textwidth}
  \includegraphics[width=\textwidth]{thumbnails/swe_phantom.png}
  \caption{SWE phantom}
\end{subfigure}
\begin{subfigure}[t]{0.24\textwidth}
  \includegraphics[width=\textwidth]{thumbnails/fi_point_scatterers.png}
  \caption{Point targets}
\end{subfigure}
\begin{subfigure}[t]{0.24\textwidth}
  \includegraphics[width=\textwidth]{thumbnails/dynamic_range_phantom.png}
  \caption{Dyn.\ range}
\end{subfigure}

\vspace{0.3em}
% Row 3
\begin{subfigure}[t]{0.24\textwidth}
  \includegraphics[width=\textwidth]{thumbnails/carotid_cross_sub2.png}
  \caption{Carotid FI}
\end{subfigure}
\begin{subfigure}[t]{0.24\textwidth}
  \includegraphics[width=\textwidth]{thumbnails/picmus_numerical.png}
  \caption{PICMUS CPWC}
\end{subfigure}
\begin{subfigure}[t]{0.24\textwidth}
  \includegraphics[width=\textwidth]{thumbnails/fieldii_stai_dynamic_range.png}
  \caption{Field~II STAI}
\end{subfigure}
\caption{Representative Delay-And-Sum (DAS) reconstructions (60\,dB dynamic range). Full catalog: \url{https://unioslo.github.io/USTB/datasets.html}}
\label{fig:previews}
\end{figure}

% ------------------------------------------------------------------
\section{Technical Approach}

\subsection*{Data Format Conversion}
Each UFF dataset is converted to OpenH-RF (\texttt{zea} HDF5) using a ready-made pipeline:
\begin{enumerate}[nosep]
  \item Read UFF with \texttt{pyuff-ustb} (channel data, probe geometry, transmit sequence).
  \item Map to \texttt{zea} scan dictionary: \texttt{probe\_geometry} [N$_\text{el}$ $\times$ 3], \texttt{polar\_angles} [N$_\text{tx}$], \texttt{focus\_distances}, \texttt{t0\_delays} [N$_\text{tx}$ $\times$ N$_\text{el}$], \texttt{sampling\_frequency}, \texttt{sound\_speed}.
  \item Generate Delay-And-Sum (DAS) B-mode image as reference reconstruction.
  \item Write \texttt{.hdf5} via \texttt{zea.File.create()}.
\end{enumerate}
\noindent Both a Python script and a MATLAB equivalent are available at \url{https://github.com/unioslo/USTB/tree/master/sandbox/OpenHRF}.

\subsection*{Reference Reconstructions \& Ground Truth}
Delay-And-Sum (DAS) beamformed images (f/1.75, Hanning apodization, coherent compounding) serve as \textbf{reference reconstructions}.
True ground truth is available where the scattering medium is fully defined:
\begin{itemize}[nosep]
  \item \textbf{Simulations}: Field~II scatterer positions and medium parameters.
  \item \textbf{Phantoms}: manufacturer-specified geometry (CIRS Model 040GSE).
  \item \textbf{SWE/ARFI}: shear-wave velocity and displacement maps.
\end{itemize}
For in-vivo data, no ground truth exists; we supply raw channel data with \textbf{processing scripts for reconstruction} as permitted by the RFP (Section~6.7).

% ------------------------------------------------------------------
\section{Data Governance, Ethics \& Compliance}

The in-vivo human datasets (cardiac and carotid) were recorded on healthy volunteers at the University of Oslo with \textbf{written informed consent} and approval from the \textbf{Regional Committee for Medical and Health Research Ethics (REK)}.
The UFF files contain only raw ultrasound channel data (backscattered RF signals) --- no identifiable patient information is stored.

Phantom and simulation data have no human-subjects implications.

% ------------------------------------------------------------------
\section{Data Rights \& Licensing}

All datasets are hosted on Zenodo under open-access terms.
Code is released under the \textbf{MIT License}.
We confirm intent to release the converted dataset under \textbf{CC~BY~4.0} as required by OpenH-RF.
No third-party IP encumbrances apply.

% ------------------------------------------------------------------
\section{Team Qualifications}

\begin{itemize}[nosep]
  \item \textbf{Ole Marius Hoel Rindal} --- Senior Lecturer at UiO and co-founder of Sonair~AS; lead developer of USTB since 2017; 15+ publications on adaptive beamforming, coherence-based imaging, and the Generalized Beamformer.
  \item \textbf{Y\"ucel Karabiyik} --- Chief Engineer and Facility Manager of the DSB ultrasound/sonar laboratory at UiO; PhD in Medical Technology (NTNU, 2017); responsible for programming and maintaining the Verasonics scanner and coordinating ultrasound data acquisitions.
  \item \textbf{Sven Peter N\"asholm} --- Professor and Head of the Digital Signal Processing and Image Analysis research group at UiO; acoustics, ultrasound signal processing, and medical imaging.
  \item \textbf{Andreas Austeng} --- Professor and Head of the Machine Learning section at UiO; 20+ years in signal processing for medical ultrasound.
\end{itemize}

% ------------------------------------------------------------------
\section{Project Plan}

\begin{table}[h!]
\centering
\small
\begin{tabular}{@{}p{5.5cm}p{3.2cm}p{4.5cm}@{}}
\toprule
\textbf{Task} & \textbf{Period} & \textbf{Deliverable} \\
\midrule
Finalize \texttt{zea} conversion pipeline & Completed & Scripts on GitHub \\
Generate reference reconstructions   & Completed & 51 DAS B-mode preview images \\
Submit proposal                      & Jun 10, 2026 & PDF via Google Form \\
Convert all datasets to OpenH-RF HDF5 & Jun 10--30, 2026 & 51 \texttt{.hdf5} files \\
Deliver data                          & Jun 30, 2026 & S3/GDrive upload \\
\bottomrule
\end{tabular}
\end{table}

\begin{thebibliography}{1}
\bibitem{rodriguez2017ustb}
A.~Rodriguez-Molares, O.~M.~H.~Rindal, O.~Bernard, A.~Nair, M.~A.~L.~Bell, H.~Liebgott, A.~Austeng, and L.~L{\o}vstakken,
``The UltraSound ToolBox,''
in \emph{Proc.\ IEEE International Ultrasonics Symposium (IUS)}, Washington, DC, USA, 2017, pp.~1--4.
\textsc{doi}: 10.1109/ULTSYM.2017.8092389.
\end{thebibliography}

\end{document}
